\begin{figure}[h!]
  \centering
  \includegraphics[width=0.5\textwidth]{figs/teaser.png}
  \caption{\textbf{Part-wise controllable 3D generation using \methodname.} Global prompts are shown at the top, and dashed arrows route reference images to target parts. Orange primitives enforce high spatial (shape-preserving) control; white primitives denote low control. \methodname localizes shape and texture edits to the designated regions.
  \vspace{-5mm}}
  \label{fig:teaser}
\end{figure}

\section{Introduction}
\label{sec:intro}

Modern 3D generative models synthesize increasingly realistic textured assets from text or image prompts~\cite{Xiang_2025_CVPR, hunyuan3d22025tencent, 10.1145/3658146}. However, their control mechanisms remain poorly aligned with practical design and editing workflows. In practice, creators rarely want a single global instruction to dictate an entire asset; rather, they require fine-grained, localized control over both shape and texture. 

For instance, a designer may require one region of an asset to strictly adhere to a target 3D structure while allowing other regions to remain flexible, leveraging the generator's learned prior to plausibly complete or refine the geometry. Similarly, they may want an appearance cue (such as a specific color, material, or reference image) to apply only to a designated component, ensuring no texture leakage occurs elsewhere. This localized intent is difficult to express with existing methods.

% Text prompts are convenient but can be geometrically ambiguous: a phrase such as ``chair with a slat back'' does not specify the exact position, dimensions, or extent of the slats. Image prompts provide stronger visual guidance, but they are difficult to edit directly in 3D and do not naturally expose local control over individual object components. Recent spatial-conditioning methods address this limitation by allowing users to guide generation with explicit 3D geometry, such as coarse primitives or proxy meshes~\cite{fedele2025spacecontrol, 10.1145/3641519.3657425, 10.1145/3641519.3657461}.

% While these methods provide a more direct interface for geometry-aware generation, they typically expose control as a single global trade-off between structural faithfulness and generative realism. Consequently, increasing the control strength preserves the input geometry but over-constrains regions that should be creatively completed, while decreasing it restores geometric plausibility at the expense of losing user-specified structure.

While text and image prompts are convenient, they are often geometrically ambiguous or difficult to edit in 3D. Recent spatial-conditioning methods address this limitation by guiding generation with explicit 3D geometry, such as coarse primitives or proxy meshes~\cite{fedele2025spacecontrol, 10.1145/3641519.3657425, 10.1145/3641519.3657461}. However, these approaches typically enforce control as a single global trade-off: increasing the control strength preserves the input geometry but over-constrains regions that should be creatively completed, while decreasing it restores geometric plausibility at the expense of losing user-specified structure.

% A similar challenge regarding locality arises in appearance control. A global text prompt like ``white sports car with green wheels'' does not explicitly bind the color ``green'' with the wheel geometry. Consequently, appearance conditions may bleed into unrelated parts of the asset or miss the target regions completely. 

% Existing appearance transfer and guidance methods improve texture adaptation by steering pretrained 3D generators or optimizing texture fields~\cite{sayandsarkar_2025_guideflow3d, 10.1145/3588432.3591503, text2mesh, Cohen-Bar_2025_CVPR}, but they are generally formulated as global or whole-object conditioning problems. Enabling controllable asset creation requires a unified interface that specifies not only \textit{what} details the model should generate, but also \textit{where} spatial and appearance controls should apply.

A similar challenge regarding locality arises in appearance control. A flat prompt like ``white sports car with green wheels'' lacks explicit spatial bindings, often causing the green texture to bleed into unrelated parts of the body or fail to cover the wheels entirely. While existing appearance transfer and guidance methods improve texture adaptation~\cite{sayandsarkar_2025_guideflow3d, 10.1145/3588432.3591503, text2mesh, Cohen-Bar_2025_CVPR}, they are generally formulated as global or whole-object conditioning problems. Enabling controllable asset creation requires a unified interface that specifies not only \textit{what} details the model should generate, but also \textit{where} spatial and appearance controls should apply.

We present \textbf{\methodname}, a training-free, test-time pipeline for part-wise controllable 3D asset generation from editable superquadric primitives. Superquadrics provide a compact and intuitive proxy for object structure: they are expressive enough to represent meaningful parts, yet simple enough to manipulate interactively. As illustrated in Fig.~\ref{fig:teaser}, each primitive can be assigned a local control level $\tau$ to govern spatial fidelity, while the same primitive layout serves as a target interface to route local text or image appearance conditions.


% In \methodname, each primitive is assigned a local control level $\tau$, indicating whether the corresponding region should adhere strictly to the input shape or adapt to the generative model's learned geometric prior. The same primitive layout also serves as an interface for part-wise text or image appearance conditions.

We evaluate \methodname on spatial-control and appearance-control tasks covering coarse object sketches, local geometric edits, local color and material conditions, and image-based appearance guidance. We compare our approach against global-control baselines and state-of-the-art appearance transfer methods. Our results show that \methodname preserves the specified geometry in high-control regions, allowing plausible variation in low-control regions and localizing appearance cues to the intended object parts without degrading overall generation quality.
 

In summary, the key contributions of our work are:
\begin{itemize}
    \item We introduce a unified, training-free pipeline for part-wise controllable 3D asset generation that uses superquadrics as a joint interface for shape and appearance.
    \item We propose a localized spatial control formulation for structure generation, enabling different regions of an asset to adopt distinct trade-offs between structural fidelity and generative realism.
    \item We design a part-aware appearance guidance scheme that routes local text or image conditions to target regions, integrated with self-similarity guidance to ensure local coherence and prevent appearance leakage.
\end{itemize}

